%% kapa: exam=1 round=27 year=2016 period=2 form=A subject=accounting subject_ko=회계학 seq=19 number=59
%% source: Q-Net 감정평가사 기출문제 (cst00309) artlSeq=5169402 file=extracted/1008/5169402/제27회 감정평가사 1차 문제/제27회 감정평가사 2교시문제(A형).pdf
%% conversion: kordoc 4.24.51 -> tools/gen1.py (auto); review=proofread-sonnet; flags=none
\begin{question}{59}
다음은 각각 독립적인 사건으로, ‘재무제표에 인식된 금액의 수정을 요하는 보고기간후사건’에 해당하는 것을 모두 고른 것은?
\begin{framedtext}
ㄱ. 보고기간말에 존재하였던 현재의무가 보고기간 후에 소송사건의 확정에 의해 확인되는 경우 \par
ㄴ. 보고기간말과 재무제표 발행승인일 사이에 투자자산의 공정가치가 하락하는 경우 \par
ㄷ. 보고기간말 이전에 구입한 자산의 취득원가나 매각한 자산의 대가를 보고기간 후에 결정하는 경우 \par
\end{framedtext}
\begin{choices}
  \item ㄱ
  \item ㄴ
  \item ㄴ, ㄷ
  \item ㄱ, ㄷ
  \item ㄱ, ㄴ, ㄷ
\end{choices}
\end{question}
\answer{4}
